
El desarrollo del sistema \textit{NodeAlert IoT} permite concluir que la transición hacia una arquitectura IoT distribuida es fundamental para garantizar la fiabilidad en entornos de seguridad crítica. A diferencia de las soluciones monolíticas, la descentralización del procesamiento en el borde (\textit{Edge Computing}) asegura que la capacidad de respuesta ante emergencias no dependa exclusivamente de la disponibilidad de la infraestructura de red global.

\begin{itemize}
    \item \textbf{Robustez del Tiempo Real:} La implementación de FreeRTOS en los nodos ESP32 ha demostrado ser un factor diferenciador. Siguiendo a \citet{valvano2014volume1}, el uso de un planificador por prioridades permite que las tareas de detección de incendio mantengan un comportamiento determinístico, eliminando las latencias críticas presentes en los sistemas basados en bucles secuenciales.
    \item \textbf{Eficiencia en la Comunicación:} El protocolo MQTT, en conjunto con la serialización JSON, ha probado ser la solución más eficiente para el intercambio de datos en redes de sensores. Como sostiene \citet{mazidi2014}, la gestión optimizada de los recursos de comunicación en microcontroladores de 32 bits permite una integración fluida entre el hardware de bajo nivel y los servicios de software de alto nivel (Django/React).
    \item \textbf{Escalabilidad e Integración:} La estructura por capas adoptada facilita la expansión del sistema. La modularidad del diseño, fundamentada en los principios de programación de \citet{deitel2004}, permite añadir nuevos nodos de sensores o diferentes tipos de actuadores sin necesidad de reestructurar el núcleo del sistema, validando la propuesta como una solución escalable para infraestructuras industriales o residenciales.
    \item \textbf{Integridad de Datos y Monitoreo:} La combinación de una base de datos centralizada con la API REST y el dashboard en React cumple con el objetivo de ubicuidad de la información. El análisis de circuitos y la estabilidad eléctrica, basados en los fundamentos de \citet{alexander2006}, garantizan que las señales capturadas sean fieles a la realidad ambiental, minimizando los falsos positivos.
\end{itemize}

En definitiva, \textit{NodeAlert IoT} representa una solución técnica integral que armoniza el diseño electrónico, la programación de sistemas embebidos y el desarrollo web, consolidando un ecosistema capaz de mitigar riesgos ambientales mediante tecnología de código abierto y alto rendimiento.